\documentclass[t,usepdftitle=false]{beamer}

\usepackage[utf8]{inputenc}
\usetheme{Singapore}
\usepackage{xcolor}
\setbeamertemplate{footline}[frame number]

% \setbeamercovered{transparent}
%\usecolortheme{crane}
\title[IFT3515]{IFT 3515\\Fonctions à plusieurs variables\\Optimisation sans contraintes\\Méthodes de région de confiance}
\author[Fabian Bastin]{Fabian Bastin\\DIRO\\Université de Montréal}
\date{}

\usepackage{enumerate}
\usepackage[francais]{babel}

\usepackage{easybmat}
\usepackage{graphicx}

\newtheorem{defn}{Définition}
\newtheorem{lem}{Lemme}
\newtheorem{thm}{Théorème}
\newtheorem{coro}{Corollaire}

\def\red{\color{red}}
\def\blue{\color{blue}}

\def\co{\mathcal{o}}

\def\cB{\mathcal{B}}
\def\cN{\mathcal{N}}
\def\cR{\mathcal{R}}

\def\bu{\boldsymbol{u}}

\newenvironment{algo}[1]{
	\refstepcounter{algo}%
	\addcontentsline{lma}{algo}{\protect\numberline{\thealgo}#1}%
	\skipbefore%
	\begin{breakbox}
		\par\noindent%
		\ifthenelse{\equal{#1}{}}{\textbf{\algoname\ \thealgo}}%
		{\textbf{\algoname\ \thealgo:} #1}
	}
	{
	\end{breakbox}
	\skipafterbox
}
\setbeamertemplate{footline}[frame number]

\begin{document}
\frame{\titlepage}

% ------------------------------------------------------------------------------------------------------------------------------------------------------

\begin{frame}
\frametitle{Contexte}

Nous considérons le problème
$$
\min_{x \in \cR^n} f(x)
$$
où $f(x) \in C^1$ et $\nabla f(x)$ est lipschitzienne sur $\cR^n$.

\mbox{}

Note: ces conditions de continuité sont parfois violées, alors que les approches continuent de donner de bons résultats.

\mbox{}

Les algorithmes d'optimisation itératifs travaillent typiquement à chaque itération avec un problème d'optimisation bien plus simple.
\begin{itemize}
\item
Dans le cas des méthodes de recherche linéaire, le sous-problème est plus simple car il est unidimensionnel.
\item
En régions de confiance, le sous-problème reste en dimension $n$, mais la fonction objectif est simplifiée et l'espace de recherche retreint.
\end{itemize}

\end{frame}

\begin{frame}
\frametitle{Sous-problème}

Conceptuellement, l'approche de régions de confiance remplace un problème à $n$ dimensions, sans contraintes, par une approximention à $n$ dimensions, soumis à des contraintes.
Avantages:
\begin{itemize}
\item
le sous-problème ne doit pas être résolu exactement;
\item
des algorithmes spécifiques existent pour traiter le sous-problème.
\end{itemize}

\begin{minipage}{0.3\textwidth}
	\includegraphics[width=\textwidth]{trbook.jpg}
\end{minipage}
\begin{minipage}{0.65\textwidth}
	La bible! A. Conn, M. Gould et Ph. L. Toint, "Trust-region methods", SIAM, 2000.
	
	\mbox{}
	
	Presque 1000 pages!
	
\end{minipage}

\end{frame}

\begin{frame}
\frametitle{Recherche linéaire vs méthodes de région de confiance}

{\red Méthodes de recherche linéaire}
\begin{enumerate}
\item
choisir une direction de descente $d_k$
\item
déterminer un longueur de pas $\alpha_k$ pour réduire suffisamment $f$ le long de $d_k$: $f(x_k + \alpha_k d_k)$.
\end{enumerate}

\mbox{}

{\red Méthodes de région de confiance}
\begin{enumerate}
\item
choisir un pas $s_k$ pour réduire le modèle de $f(x_k+s)$
\item
accepter $x_{k+1} = x_k + s_k$ si la réduction prédite par le modèle est vérifiée en $f(x_k + s_k)$
\item
autrement poser $x_{k+1} = x_k$ et raffiner le modèle.
\end{enumerate}

\end{frame}

\begin{frame}
\frametitle{Modèle}

Nous voudrions un modèle $m_k$ de $f$ en $x_k$ qui satisfait au minimum les hypothèses suivantes:
\begin{align*}
& m_k \in C^1 \\
& m_k(0) = f(x_k) \\
& \nabla m_k(0) = \nabla f(x_k)
\end{align*}

\mbox{}

Généralement, un modèle quadratique est sélectionné:
$$
m_k(s) = f(x_k) + \nabla f(x_k)^Ts + \frac{1}{2}s^TH_ks
$$
avec $H_k$ symétrique (mais pas nécessairement définie positive!).

\end{frame}

\begin{frame}
\frametitle{Région de confiance}

Pour de petits $s$, le modèle devrait ressembler à la fonction $f$, mais plus $s$ est grand, plus il est probable que $m$ et $f$ soient différents!

\mbox{}

Nous allons contrer ce problème en bornant la norme du pas:
$$
\| s_k \|_k \leq \Delta_k
$$
où $\Delta_k$ est le {\red rayon de la région de confiance} et $\| \cdot \|_k$ est une norme, dont le choix peut être dépendant de l'itération $k$.
Le plus souvent toutefois, nous travaillerons avec la norme 2: $\| \cdot \|_2$.

\mbox{}

Dès lors, le sous-problème de région de confiance est le problème sous contraintes
\begin{align*}
\min_s \ & m_k(s) \\
\mbox{t.q. } & \| s \|_k \leq \Delta_k.
\end{align*}

\end{frame}

\begin{frame}
\frametitle{Région de confiance}

La {\red région de confiance} à l'itération $k$ est définie comme
\[
\mathcal{B}_k = \left\lbrace x \in \cR^n \,|\, \| x - x_k \|_k \leq
\Delta_k \right\rbrace,
\]

\mbox{}

Outre $\| \cdot\|_2 $, des choix usuels pour la norme sont $\| \cdot \|_1$ et $\| \cdot \|_{\infty}$.

\mbox{}

Rappel: si $x = (x_1, x_2,\ldots, x_n)$,
\begin{align*}
\| x \|_2 & = \sqrt{x^T x} \\
\| x \|_1 & = \sum_{i = 1}^n | x_i | \\
\| x \|_\infty & = \max \{ |x_i|,\ i = 1,\ldots,n \}
\end{align*}

\end{frame}

\begin{frame}
\frametitle{Modèle quadratique}

Si $f \in C^2$, on pourra choisir pour $H_k$ la matrice hessienne $\nabla^2 f(x_k)$, mais comme dans le cadre de la recherche linéaire, elle peut être coûteuse à obtenir.

\mbox{}

Si $f \in C^1$ et son gradient est une fonction lipschitzienne, on pourra recourir à des approximations, comme l'approche BFGS. Mais comme le caractère défini positif n'est plus requis, d'autres options s'ouvrent, en particulier l'approximation SR1.

\end{frame}

\begin{frame}
\frametitle{Algorithme basique de région de confiance}

%\begin{algo}{Basic trust-region algorithm (BTR)}
\begin{description}
\item[Étape 0. Initialisation]
Un point initial $x_0$ et un rayon initial de région de confiance $\Delta_0$ sont donnés.
Les constantes $\eta_1$, $\eta_2$, $\gamma_1$ et $\gamma_2$ sont aussi données et satisfont
$$
0 < \eta_1 \leq \eta_2 < 1 \text{ et } 0 < \gamma_1 \leq \gamma_2 < 1.
$$
Calculer $f(x_0)$ et poser $k = 0$.

\item[Étape 1. Définition du modèle]
Choisir $\|\cdot\|_k$ et définir un modèle $m_k$ dans $\mathcal{B}_k$.

\item[Étape 2. Calcul du pas]

Calculer un pas $s_k$ qui ``réduit suffisamment le modèle'' $m_k$ et tel que $x_k + s_k \in \mathcal{B}_k$.

\item[Étape 3. Acceptation du point candidat]
Calculer $f(x_k+s_k)$ et définir
\begin{equation}
\rho_k = \frac{f(x_k)-f(x_k+s_k)}{m_k(0)-m_k(s_k)}.
\label{eq:rhok}
\end{equation}
Si $\rho_k \geq \eta_1$, définir $x_{k+1} = x_k+s_k$; autrement définir $x_{k+1} = x_k$.
\end{description}

\end{frame}

\begin{frame}
\frametitle{Algorithme basique de région de confiance}

\begin{description}
\item[Étape 4. Mise à jour de la région de confiance]
		\[
		\Delta_{k+1} \in \begin{cases}
		[\Delta_k, \infty) & \text{si } \rho_k \geq \eta_2,\\
		[\gamma_2\Delta_k, \Delta_k] & \text{si } \rho_k \in [ \eta_1, \eta_2), \\
		[\gamma_1\Delta_k, \gamma_2\Delta_k] & \text{si } \rho_k < \eta_1.
		\end{cases}
		\]
Incrémenter $k$ de 1 et retourner à l'étape 1.
\end{description}
%\end{algo}

\mbox{}

Dans cette description, des valeurs raisonnables pour les constantes sont, par exemple,
\[
\eta_1 = 0.01,\ \eta_2 = 0.9,\ \text{ et } \gamma_1 = \gamma_2 = 0.5,
\]
mais d'autres valeurs peuvent être sélectionnées.

\end{frame}

\begin{frame}
\frametitle{Itérations}

\begin{description}
\item[Itération réussie]
Si $\rho_k \geq \eta_1$ à l'étape 1, l'itération~$k$ est dite réussie puisque le point candidat $x_k + s_k$ est accepté.
\item[Itération non réussie]
Sinon, l'itération est déclarée non réussie et le point candidat est rejeté. 
\item[Itération très réussie]
Si $\rho_k \geq \eta_2$, l'accord entre le modèle et la fonction est particulièrement bon, aussi l'itération est dite très réussie. Ceci suggère d'élargir la région de confiance comme décrit à l'étape 4, afin de permettre un plus long pas à la prochaine itération.
\end{description}
\end{frame}

\begin{frame}
\frametitle{Solution approximative du sous-problème}

Chaque sous-problème de région de confiance doit être résolu approximativement, et aussi économiquement que possible.

\mbox{}

Afin d'être capable de garantir la convergence globale de la méthode, nous voudrions au minimum que la solution approximative atteigne le même niveau de réduction de la valeur du modèle que ne le ferait un pas dans la direction de plus forte pente, contraint par la région de confiance.

\end{frame}

\begin{frame}
\frametitle{Pas de Cauchy}

Le pas de Cauchy est défini par $s_k^C := -\alpha_k^C \nabla f(x_k)$, où
\begin{align*}
\alpha_k^C :&= \arg\min
\left\{ m_k(-\alpha \nabla f(x_k)) \,|\, \alpha > 0, \alpha\| \nabla f(x_k) \| \leq \Delta_k \right\} \\
&= \arg\min \left\{ m_k(x-\alpha \nabla f(x_k)) \,|\, 0 < \alpha \leq  \Delta_k/\| \nabla f(x_k) \| \right\}
\end{align*}

\mbox{}

Calculer le pas de Cauchy est très facile, comme il s'agit de la minimisation d'une fonction quadratique sur un segment de droite.

\mbox{}

Habituellement, nous exigerons que la solution approximative du sous-problème de région de confiance entraîne une réduction du modèle au moins équivalente à celle obtenue avec le pas de Cauchy:
$$
m_k(s_k) \leq m_k(s_k^C),\ \| s_k \|_k \leq \Delta_k.
$$

\end{frame}

\begin{frame}
\frametitle{Calcul du pas de Cauchy}

Nous devons minimiser
$$
m_k(-\alpha \nabla f(x_k)) = f(x_k) - \alpha \| \nabla f(x_k) \|_2^2 + \frac{1}{2} \alpha^2 \nabla f(x_k)^T H_k \nabla f(x_k)
$$
dont la dérivée en $\alpha$ est
$$
\frac{d m_k(-\alpha \nabla f(x_k))}{d \alpha} = -\| \nabla f(x_k) \|_2^2 + \alpha 
\nabla f(x_k)^T H_k \nabla f(x_k)
$$
Supposons $m_k(\cdot)$ strictement convexe en $\alpha$, i.e. $\nabla f(x_k)^T H_k \nabla f(x_k) > 0$.
Le minimum de $m_k(\cdot)$ le long de $-\alpha \nabla f(x_k)$ est obtenu en annulant la dérivée par rapport à $\alpha$:
$$
\alpha^* = \frac{\| \nabla f(x_k) \|_2^2}{\nabla f(x_k)^T H_k \nabla f(x_k)}
$$
Mais il se peut que ce minimum soit hors de la région de confiance, i.e.
$$
\alpha^* \| \nabla f(x_k) \|_k \geq \Delta_k
$$

\end{frame}

\begin{frame}
\frametitle{Calcul du pas de Cauchy (suite)}

Dès lors, pour assurer de rester de rester dans la région de confiance, nous prendrons
$$
\alpha^* = \min
\left\{
\frac{\Delta_k}{\| \nabla f(x_k) \|_k},
\frac{\| \nabla f(x_k) \|_2^2}{\nabla f(x_k)^T H_k \nabla f(x_k)}
\right\}
$$

\mbox{}

Si $m_k(\cdot)$ n'est pas strictement convexe en $\alpha$, le modèle n'est pas borné inférieurement le long de $-\alpha \nabla f(x_k))$, et il faut s'arrêter sur la frontière de la région de confiance, ce qui donne
$$
\alpha^* = \frac{\Delta_k}{\| \nabla f(x_k) \|_k}
$$
\end{frame}

\begin{frame}
\frametitle{Théorie de convergence}

\begin{thm}[Décroissance du modèle atteignable]
Soit $m_k(s)$ un modèle quadratique de $f$ autour de $x_k$, et soit $s_k^C$ le pas de Cauchy, satisfaisant $\| s_k^C \|_k \leq \Delta_k$.
Alors la diminution atteignable du modèle est d'au moins
$$
f(x_k) - m_k(s_k^C) \geq \frac{1}{2} \| \nabla f(x_k) \|_2
\min \left\{ \frac{\| \nabla f(x_k) \|}{1+\| H_k \|_2}, \kappa_s \Delta_k \right\}
$$
\end{thm}

Note: $\kappa_s$ est un facteur d'équivalence de normes. On choisit la norme telle que
$$
\kappa_s \| \cdot \|_k \leq  \| \cdot \|_2 \leq \kappa_l \| \cdot \|_k
$$
avec $\kappa_l \geq \kappa_s > 0$. Les normes usuelles satisfont cette condition:
\begin{align*}
\frac{1}{\sqrt{n}} \| \cdot \|_1 & \leq \| \cdot \|_2 \leq \| \cdot \|_1 \\
 \| \cdot \|_{\infty} & \leq \| \cdot \|_2 \leq n \| \cdot \|_{\infty} \\
\end{align*}

\end{frame}

\begin{frame}
\frametitle{Théorie de convergence}

\begin{coro}
Soit $m_k(s)$ un modèle quadratique de $f$ et $s_k$ tel que $m_k(s_k) \leq m_k(s_k^C)$, et $\| s_k \|_k \leq \Delta_k$.
Alors
$$
f(x_k) - m_k(s_k) \geq \frac{1}{2} \| \nabla f(x_k) \|_2
\min \left\{ \frac{\| \nabla f(x_k) \|}{1+\| H_k \|_2}, \kappa_s \Delta_k \right\}
$$
De plus, si l'itération $k$ est très réussie,
$$
f(x_k) - f(x_{k+1}) \geq \eta_2 \frac{1}{2} \| \nabla f(x_k) \|_2
\min \left\{ \frac{\| \nabla f(x_k) \|}{1+\| H_k \|_2}, \kappa_s \Delta_k \right\}
$$
\end{coro}

\end{frame}

\begin{frame}
\frametitle{Adéquation du modèle}

\begin{lem}[Différence entre le modèle et la fonction]
Soit $f \in C^2$, et supposons qu'il existe des constantes $\kappa_H \geq 1$ et $\kappa_m \geq 0$ telles que $\| \nabla^2 f(x_k) \|_2 \leq \kappa_H$ et $\| H_k \|_2 \leq \kappa_m$ pour tout $k$. Alors
$$
| f(x_k+s_k) - m_k(s_k) | \leq \kappa_d \Delta_k^2
$$
où $\kappa_d = \frac{1}{2} \kappa_l^2 (\kappa_H + \kappa_m)$.
\end{lem}

\end{frame}

\begin{frame}
\frametitle{Comportement du rayon de la région de confiance}

L'algorithme garantit que le rayon de la région de confiance ne peut être arbitrairement petit à des points non optimums.
% converge pas vers 0 quand l'indice d'itération tend vers l'infini.

\begin{thm}
	Soit $f \in C^2$, et supposons qu'il existe des constantes $\kappa_H \geq 1$ et $\kappa_m \geq 0$ telles que $\| \nabla^2 f(x_k) \|_2 \leq \kappa_H$ et $\| H_k \|_2 \leq \kappa_m$ pour tout $k$.
	Soit $\kappa_d = \frac{1}{2} \kappa_l^2 (\kappa_H + \kappa_m)$.
	Si à l'itération $k$, nous avons $\nabla f(x_k) \ne 0$ et
	$$
	\Delta_k \leq \| \nabla f(x_k) \|_2
	\min \left\{ \frac{1}{\kappa_s(\kappa_H + \kappa_b)}, \frac{\kappa_s(1-\eta_2)}{2\kappa_d} \right\},\ \forall k,
	$$
	alors l'itération $k$ est très réussie et $\Delta_{k+1} \geq \Delta_k$.
\end{thm}

\end{frame}

\begin{frame}
\frametitle{Comportement du rayon de la région de confiance}

\begin{coro}
Soit $f \in C^2$, et supposons qu'il existe des constantes $\kappa_H \geq 1$ et $\kappa_m \geq 0$ telles que $\| \nabla^2 f(x_k) \|_2 \leq \kappa_H$ et $\| H_k \|_2 \leq \kappa_m$ pour tout $k$.
Soit $\kappa_d = \frac{1}{2} \kappa_l^2 (\kappa_H + \kappa_m)$.
S'il existe une constante $\epsilon > 0$ telle que $\| \nabla f(x_k) \|_2 \geq \epsilon$ pour tout $k$, alors
$$
\Delta_k \geq \kappa_{\epsilon} := \epsilon \gamma_1
\min \left\{ \frac{1}{\kappa_s(\kappa_H + \kappa_b)}, \frac{\kappa_s(1-\eta_2)}{2\kappa_d} \right\},\ \forall k.
$$
\end{coro}

\begin{coro}[Possible arrêt en un nombre fini d'itérations]
Soit $f \in C^2$, et supposons que les matrices hessiennes de l'objectif et du modèle uniformément bornées positivement.
Si la méthode de région de confiance a seulement un nombre fini d'itérations réussies, alors $x_k = x^*$ et $\nabla f(x^*)$ pour tout $k$ assez grand.
\end{coro}

\end{frame}

\begin{frame}
\frametitle{Convergence globale}

\begin{thm}[Convergence globale]
Soit $f \in C^2$, et supposons que les matrices hessiennes de l'objectif et du modèle uniformément bornées positivement.
Un des trois cas suivants a lieu:
\begin{enumerate}
\item
$\nabla f(x_k) = 0$ pour un certain $k \in \cN$
\item
$\lim_{k \rightarrow \infty} f(x_k) = -\infty$
\item
$\lim_{k \rightarrow \infty} \nabla f(x_k) = 0$
\end{enumerate}
\end{thm}

\end{frame}

\begin{frame}
\frametitle{Méthodes pour résoudre le sous-problème de région de confiance}

Examinons plus en détails comment résoudre le sous-problème de région de confiance
\begin{align*}
\min_{s \in \cR^n}\ & m(s) = s^T \nabla f(x_k) +  \frac{1}{2} s^T \nabla^2 f(x_k) s \\
\mbox{t.q. } & \| s \| \leq  \Delta_k.
\end{align*}
de telle sorte que la théorie de convergence ci-dessus s'applique.

\mbox{}

En d'autres mots, nous cherchons $s^* \in \cR^n$ tel que
$$
m(s^*) \leq m(s^c) \mbox{ et } \| s^* \| \leq \Delta.
$$
On pourrait résoudre
\begin{itemize}
\item
exactement $\rightarrow$ méthode de type Newton
\item
approximativement
\end{itemize}

\end{frame}

\begin{frame}
\frametitle{Sous-problème}

À partir de maintenant, nous utiliserons la norme 2 pour déterminer les régions de confiance, aussi avons-nous à résoudre approximativement
$$
(TRS)\ \min_{s \in \cR^n} m(s) \mbox{ s.t. } \| s \|_2 \leq \Delta
$$
où $m(s) = \nabla f(x_k)^Ts + \frac{1}{2} s^T H s$.

\begin{thm}
N'importe quel minimiseur global $s^*$ de $(TRS)$ doit satisfaire
\begin{enumerate}
\item 
$(H + \lambda^* I)s^* = - \nabla f(x_k)$,
\item
$H + \lambda^* I$ semi-définie positive
\item
$\lambda^* \geq 0$
\item
$ \lambda^*(\| s^* \|_2 - \Delta ) = 0$.
\end{enumerate}
De plus, si $H + \lambda^* I$ est définie positive, alors $s^*$ est unique.
\end{thm}

\end{frame}

\begin{frame}
\frametitle{Solutions exactes de (TRS)}

Si $H$ est définie positive et la solution de $Hs = -\nabla f(x_k)$ satisfait $\| s \|_2 \leq \Delta$, alors $s^*  = s$.

\mbox{}

Si $H$ est indéfinie ou si la solution du système $Hs = -\nabla f(x_k)$ satisfait $\| s \|_2 > \Delta$, résoudre le système nonlinéaire
$$
(H + \lambda I)s = -\nabla f(x_k),\mbox{ t.q. } \| s \| = \Delta,
$$
pour $s$ et $\lambda$ en utilisant la méthode de Newton.
Des difficultés surviennent
\begin{itemize}
\item
possiblement quand de multiples solutions locales apparaissent,
\item
ou quand $\nabla f(x_k)$ est proche de l'orthogonalité avec le(s) vecteur(s) propre(s) correspondant à la valeur propre la plus négative de $H$.
\end{itemize}
Quand $n$ est grand, la factorisation pour résoudre $Hs = -g$ peut être impossible.

\end{frame}

\begin{frame}
\frametitle{Solutions approximatives de (TRS)}

Comme nous n'avons besoin que d'une solution approximative, il est préférable d'employer une méthode itérative.
\begin{enumerate}
\item
La méthode de plus forte pente mène au pas de Cauchy $s^C$.
\item
Utiliser la méthode des gradients conjugués pour améliorer le pas à partir de $s^C$. Point à considérer:
\begin{itemize}
	\item il faut rester dans la région de confiance,
	\item comment gérer une courbure négative?
\end{itemize}
\end{enumerate}

\end{frame}

\end{document}